\section{限制、工程判断与未来方向}

课程最后的 summary 与 future work 很克制：贡献是一个简单、稳定、可扩展的新视角，而不是宣布 RL 已经被 language modeling 取代。本节把课堂 Q\&A 中最重要的限制整理为可执行检查项。

\subsection{课程总结：方法为什么值得重视}

Decision Transformer 的价值来自三点。第一，用 conditional sequence modeling 统一 return、state 与 action；第二，用稳定监督目标训练，不显式依赖 value function 或 policy gradient；第三，在多类 offline benchmark 上达到有竞争力的结果，并通过 context、sparse reward 与 critic 实验展示 sequence history 的作用。

\lecturefigure{slide-22-summary.jpg}{原讲座总结：简单、稳定、易训练，并在多类 RL 设置中表现强劲。}{Stanford Online 官方视频 1:08:56。}

\begin{importantbox}{证据比例原则}
可以说：Decision Transformer 证明了 sequence modeling 是可行且强的 offline-RL formulation。不能说：它证明所有 RL 都应被 next-token prediction 取代，或只要扩大 Transformer 就能解决 exploration、safety 和 distribution shift。
\end{importantbox}

\subsection{Online RL 不能靠简单加噪声获得}

在 offline setting 中，训练数据已经存在；online RL 必须一边收集数据一边改进 policy。课堂 Q\&A 询问能否直接采用 epsilon-greedy 或 Boltzmann sampling。讲者的回答是：初步证据表明传统 RL 的 exploration recipe 不能原封不动迁移，exploration/exploitation balance 与错误轨迹如何进入 context 都需要新的设计。

\begin{warningbox}{在线化的三个额外问题}
\begin{enumerate}
  \item \textbf{Cold start}：没有高回报 trajectory 时，target conditioning 缺少对应行为；
  \item \textbf{Non-stationary dataset}：policy 更新会持续改变数据分布；
  \item \textbf{Sequence contamination}：探索动作产生的低质量历史会影响后续预测，不能简单视为 iid augmentation。
\end{enumerate}
\end{warningbox}

\teachervoice{讲义提醒：讲者没有声称 epsilon-greedy 必然失败，而是说“devil lies in the detail”，现成在线 RL 技巧并非直接 transferable。正确研究问题是如何设计数据收集、目标条件与 sequence replay，而不是只调一个 sampling temperature。}

\subsection{Discount factor 与有限 context}

原实验主要使用 $\gamma=1$ 的 finite-horizon return，但 Decision Transformer 与 discounting 并不冲突。只需在数据预处理和 rollout 中一致地使用 $\widehat{R}^{(\gamma)}_t$。然而 discounting 解决的是远期 reward 权重，不等于 memory：即使 $\gamma$ 合理，有限 $K$ 仍可能看不到完成任务所需的早期事件。

\begin{importantbox}{Discount 与 context 回答不同问题}
$\gamma$ 决定远期 reward 的数值权重；$K$ 决定模型可访问多少历史信息。较小 $\gamma$ 不能恢复被截断的关键 observation，较大 $K$ 也不能替代对无限时域 return 的稳定定义。
\end{importantbox}

\subsection{与 CQL、pessimism 和 hybrid objective 的关系}

CQL（Conservative Q-Learning）通过压低数据外 action 的估计价值来对抗 offline overestimation。课堂 Q\&A 提出能否把 pessimism 或 Q-learning objective 加到 Decision Transformer 上。答案是原则上可以，但原工作的科学问题恰恰是：不加入这些传统组件，仅用简单 sequence objective 能走多远。

\begin{knowledgebox}{三类 hybrid 方向}
\begin{enumerate}
  \item \textbf{Sequence actor + pessimistic critic}：DT 生成候选动作，保守 value model 过滤；
  \item \textbf{Return model + uncertainty}：对 desired return 的可达性与置信度联合建模；
  \item \textbf{Online sequence replay}：按 trajectory quality、novelty 与 safety 维护不断更新的上下文数据。
\end{enumerate}
这些方向增加了算法复杂度，也要求重新验证稳定性与因果贡献。
\end{knowledgebox}

\subsection{未来工作：multimodal、multitask、multi-agent}

讲者的 future-work slide 把序列建模视为统一接口：multimodality 让模型同时处理在线与离线 experience、图像和文本；multitask 需要更丰富的 task-conditioned goal 或语言；multi-agent 则把其他 agent 的动作和 observation 纳入联合序列。每个方向都扩大 token space，也同步扩大 credit、coordination 与 safety 难题。

\lecturefigure{slide-23-future-work.jpg}{未来方向：multimodality、multitask 与 multi-agent sequential decision making。}{Stanford Online 官方视频 1:10:46。}

\begin{knowledgebox}{读图：统一序列接口不等于统一难度}
把所有模态转成 token 可以统一软件栈，却不会自动统一 reward semantics、action latency、agent identity 或安全约束。Multimodal model 需要时间同步，multitask model 需要避免目标混淆，multi-agent model 还要处理非平稳对手与联合 credit assignment。
\end{knowledgebox}

\subsection{原讲座资源入口}

原 slide 给出项目页、论文和代码。它们分别承担不同角色：项目页提供可视化与简要说明，论文给出方法和实验定义，代码揭示数据预处理、sequence sampling、return scale 与评估细节。复现时应从代码和原始数据配置核对，而不是只根据架构图重写。

\lecturefigure{slide-24-useful-links.jpg}{原讲座列出的 Decision Transformer 项目页、论文与代码。}{Stanford Online 官方视频 1:11:30。}

\begin{importantbox}{复现检查清单}
\begin{enumerate}
  \item 核对 reward normalization、return scale 与 target-return 选择；
  \item 记录 dataset collection policy、trajectory length 与成功比例；
  \item 对齐 context length、sequence sampling 与 padding mask；
  \item 同时报训练 loss、realized return、target sweep 与多随机种子方差；
  \item 至少比较 BC、percent BC、一个 value-based offline baseline 与 context ablation。
\end{enumerate}
\end{importantbox}

\subsection{本章小结}

Decision Transformer 的边界与贡献同样重要。它提供稳定而统一的 offline sequence objective，但 online exploration、超出数据支持的行为、有限 context、discounting 与 conservative policy improvement 仍需单独解决。

